%==========================================================
%  CHAPITRE 2 — CONTEXTE GÉNÉRAL DU PROJET
%==========================================================
\chapter{Contexte général du projet}
\label{chap:contexte}

%──────────────────────────────────────────────────────────
\section{Étude de l'existant}
\label{sec:ctx:existant}
%──────────────────────────────────────────────────────────

\subsection{Situation actuelle}

\todo{Décrire le système ou le processus existant avant le projet :
outils utilisés, fonctionnement manuel ou semi-automatisé, acteurs
impliqués, flux d'information.}

\todo{Si applicable, insérer un schéma du processus existant (diagramme BPMN,
flowchart) avec \textbackslash placeholder\{...\}.}

\placeholder{Schéma du processus existant (ex : diagramme BPMN, flowchart —
à insérer dans figures/existant.png)}

\subsection{Limites et problèmes de l'existant}

\begin{center}
\begin{tabular}{>{\sffamily\bfseries}p{4cm} p{9cm}}
  \toprule
  \textbf{Problème identifié} & \textbf{Impact} \\
  \midrule
  \todo{Saisie manuelle}
    & \todo{Risque d'erreurs, perte de temps (ex : 2h/jour)} \\
  \todo{Absence de traçabilité}
    & \todo{Impossible de retrouver l'historique des opérations} \\
  \todo{Silos de données}
    & \todo{Données dispersées dans plusieurs fichiers Excel non synchronisés} \\
  \todo{Scalabilité limitée}
    & \todo{Le système actuel ne supporte pas une charge > N utilisateurs simultanés} \\
  \bottomrule
\end{tabular}
\end{center}

%──────────────────────────────────────────────────────────
\section{Analyse des besoins}
\label{sec:ctx:besoins}
%──────────────────────────────────────────────────────────

\subsection{Identification des acteurs}

\begin{center}
\begin{tabular}{>{\sffamily\bfseries}l l p{7cm}}
  \toprule
  \textbf{Acteur} & \textbf{Type} & \textbf{Rôle dans le système} \\
  \midrule
  \todo{Administrateur} & Interne & \todo{Gestion des utilisateurs et des droits} \\
  \todo{Gestionnaire}   & Interne & \todo{Saisie et validation des données métier} \\
  \todo{Utilisateur}    & Interne & \todo{Consultation des tableaux de bord} \\
  \todo{Client / API tierce} & Externe & \todo{Déclenchement d'actions via API} \\
  \bottomrule
\end{tabular}
\end{center}

\subsection{Besoins fonctionnels}

\begin{enumerate}
  \item \textbf{BF-01 :} \todo{Description du besoin fonctionnel 1.
        Exemple : Le système doit permettre à un gestionnaire de créer,
        modifier et supprimer une commande.}
  \item \textbf{BF-02 :} \todo{Description du besoin fonctionnel 2.}
  \item \textbf{BF-03 :} \todo{Description du besoin fonctionnel 3.}
  \item \textbf{BF-04 :} \todo{Description du besoin fonctionnel 4.}
  \item \textbf{BF-05 :} \todo{Description du besoin fonctionnel 5.}
  \item \todo{Ajouter d'autres besoins fonctionnels.}
\end{enumerate}

\subsection{Besoins non fonctionnels}

\begin{itemize}
  \item \textbf{Performance :} \todo{ex : temps de réponse < 300 ms pour 95\% des requêtes}
  \item \textbf{Disponibilité :} \todo{ex : SLA 99,5\% (hors maintenance planifiée)}
  \item \textbf{Sécurité :} \todo{ex : authentification JWT, chiffrement TLS 1.3, RBAC}
  \item \textbf{Maintenabilité :} \todo{ex : couverture de tests > 75\%, documentation Swagger}
  \item \textbf{Scalabilité :} \todo{ex : architecture permettant la montée en charge horizontale}
  \item \textbf{Ergonomie :} \todo{ex : interface responsive, accessible WCAG 2.1 niveau AA}
\end{itemize}

%──────────────────────────────────────────────────────────
\section{Cahier des charges}
\label{sec:ctx:cdc}
%──────────────────────────────────────────────────────────

\subsection{Périmètre du projet}

\todo{Délimiter précisément ce qui est dans le périmètre (in-scope)
et hors périmètre (out-of-scope) du projet.}

\begin{center}
\begin{tabularx}{\textwidth}{X X}
  \toprule
  \rowcolor{ChapBlue!10}
  \textbf{Dans le périmètre (In-scope)} & \textbf{Hors périmètre (Out-of-scope)} \\
  \midrule
  \todo{Fonctionnalité A} & \todo{Module existant non touché} \\
  \todo{Fonctionnalité B} & \todo{Intégration ERP tiers} \\
  \todo{Fonctionnalité C} & \todo{Application mobile} \\
  \bottomrule
\end{tabularx}
\end{center}

\subsection{Contraintes du projet}

\begin{itemize}
  \item \textbf{Contraintes techniques :} \todo{stack imposée, compatibilité
        avec l'infrastructure existante, langage imposé par l'équipe}
  \item \textbf{Contraintes temporelles :} \todo{livraison V1 attendue à J+6 semaines}
  \item \textbf{Contraintes budgétaires :} \todo{utilisation de solutions open-source uniquement}
  \item \textbf{Contraintes réglementaires :} \todo{conformité RGPD, ISO 27001, etc.}
\end{itemize}

%──────────────────────────────────────────────────────────
\section{Solutions étudiées et choix technologiques}
\label{sec:ctx:solutions}
%──────────────────────────────────────────────────────────

\todo{Si plusieurs solutions ont été étudiées avant de choisir la
solution finale, présenter ici le benchmark ou la matrice de décision.}

\begin{center}
\begin{tabular}{>{\sffamily\bfseries}l c c c c c}
  \toprule
  \rowcolor{ChapBlue!10}
  \textbf{Critère / Solution} & \textbf{Poids}
    & \textbf{\todo{Sol. A}} & \textbf{\todo{Sol. B}} & \textbf{\todo{Sol. C}}
    & \textbf{Retenu} \\
  \midrule
  Performance   & 30\% & 4/5 & 3/5 & 5/5 & \todo{Sol. C} \\
  Coût          & 20\% & 5/5 & 3/5 & 4/5 & \\
  Maintenabilité& 25\% & 3/5 & 5/5 & 4/5 & \\
  Communauté    & 15\% & 4/5 & 4/5 & 5/5 & \\
  Intégration   & 10\% & 5/5 & 2/5 & 4/5 & \\
  \midrule
  \textbf{Score pondéré} & & \todo{XX} & \todo{XX} & \todo{XX} & \\
  \bottomrule
\end{tabular}
\end{center}

\clearpage
